% STR -- os/samir/core/fn_builtins.c fn_str (l.361-392)
\begin{XbaseEntry}{STR}{function}
Converts a number to a character string, as a character string.
\Syntax
\SyntaxLine{\kw{STR(}\xp{expN1} \opt{\kw{,} \xp{expN2} \opt{\kw{,} \xp{expN3}}}\kw{)}}
\Defaults
\xp{expN2}, the length of the result, is 10. \xp{expN3}, the number of decimal places, is 0. SET DECIMALS has no effect on STR(); the places are those of \xp{expN3} only.
\Usage
The result always has exactly \xp{expN2} characters. The number is rounded to \xp{expN3} places, with a value half-way between two results rounded toward the larger number, and it is placed at the right of the field with blanks on the left. A negative sign and the decimal point each take one character of the length.

When the number does not fit in the length, every character of the result is an asterisk. The length must be from 1 through 254, the number of places from 0 through 15, and, when the places are not 0, the length must be at least the places plus 2. Any other value raises \xmsg{63}{STR() : Out of range.}. A non-numeric argument raises \xmsg{9}{Data type mismatch.}, and a call with no argument or with more than three raises \xmsg{11}{Invalid function argument.}.
\Tips
The result of \fn{STR()} with its default length begins with blanks. Apply \fn{LTRIM()} to remove them before joining the number to text.
\Examples
\DotExample{examples/fn-str}
A number too wide for its field fills the field with asterisks; a length too small for the places is refused:
\DotExample{examples/fn-str-2}
\SeeAlso{VAL(), LTRIM(), ROUND(), INT(), TRANSFORM()}
\end{XbaseEntry}

